\chapter{Record and Check Demonstrations}
\label{chap:data}

The MuJoCo teleoperation program records complete task attempts as episodes in a
Zarr dataset. Recording and control occur in the same program, so the saved
images, measured state, and demonstrated target share one timeline.

\section{Record Episodes}

Start \filepath{./launch_all.sh} as described in \cref{chap:touch}. The MuJoCo
window uses these controls:

\begin{table}[H]
  \centering
  \begin{tabular}{@{}ll@{}}
    \toprule
    \textbf{Key} & \textbf{Action} \\
    \midrule
    \filepath{Space} & Start recording; press again to stop and save the episode. \\
    \filepath{Backspace} & Cancel the current episode without saving it. \\
    \filepath{R} & Reset the robot and randomize the scene. \\
    \filepath{Esc} & Quit. \\
    \bottomrule
  \end{tabular}
\end{table}

The recorder creates a timestamped dataset under \filepath{data/datasets/} and
samples at 10~Hz. Practice before recording. Start just before useful motion,
finish after task completion, and cancel failed, interrupted, or corrupted
attempts rather than teaching them as intended behavior.

\section{Define one synchronized sample}

A visuomotor sample normally contains:

\begin{equation}
  \left(o_t, a_t\right)
  = \left(\text{images}_t,\text{robot state}_t,
  \text{demonstrated command}_t\right).
\end{equation}

An episode is the ordered sequence from one start/stop interval. Episode
boundaries must remain intact because a training window must never cross from
the end of one attempt into the reset of the next.

\section{Define and document the schema}

The exact keys are project-specific. For example, the internship recorder
stored the following Zarr arrays at 10 Hz:

\begin{table}[H]
  \centering
    \caption{MuJoCo dataset interface used by the project recorder.}
  \begin{tabularx}{\textwidth}{@{}p{0.34\textwidth}p{0.20\textwidth}X@{}}
    \toprule
    \textbf{Key} & \textbf{Per-step shape} & \textbf{Meaning} \\
    \midrule
    \filepath{agentview_image} & $84\times84\times3$ & Global RGB image. \\
    \filepath{robot0_eye_in_hand_image} & $84\times84\times3$ & Wrist RGB image. \\
    \filepath{robot0_eef_pos} & 3 & Measured end-effector position. \\
    \filepath{robot0_eef_quat} & 4 & Measured end-effector quaternion. \\
    \filepath{robot0_gripper_qpos} & 1 & Measured gripper opening/state. \\
    \filepath{action} & 8 & Target pose and gripper command. \\
    \bottomrule
  \end{tabularx}
\end{table}

The simulation action is

\begin{equation}
  a_t=[x,y,z,q_x,q_y,q_z,q_w,g].
\end{equation}

The quaternion stored in the action is \filepath{xyzw}; the measured
\filepath{robot0_eef_quat} observation is \filepath{wxyz}. This project-specific
difference is handled by the current recorder and runner, but must not be
silently changed or mixed with real-robot datasets.

\section{Inspect structure, content, and coverage}

Validate every new dataset before sending it to the training server.

First confirm episode count, equal array lengths, action dimension, image
shapes, and valid Zarr structure:

\begin{lstlisting}[language=bash]
python scripts/check_zarr.py data/datasets/<dataset>.zarr
\end{lstlisting}

Then replay every episode and inspect the robot trajectory:

\begin{lstlisting}[language=bash]
python scripts/replay_zarr.py data/datasets/<dataset>.zarr --scale 2
python scripts/traj_zarr.py data/datasets/<dataset>.zarr
\end{lstlisting}

Look for frozen or repeated images, long idle intervals, abrupt target jumps,
unintended corrective loops, incorrect gripper transitions, and actions that
cannot be explained from the recorded observations. Also verify that object
positions and orientations cover the range expected during policy execution.

\begin{methodbox}[Demonstration quality]
A successful episode is not automatically a good demonstration. Keep motions
deliberate and consistent, avoid hiding the object at critical moments, and
cancel an episode when the operator makes a correction that should not be
imitated. Record three episodes first and inspect all of them before continuing
with a long session.
\end{methodbox}

Keep the original dataset immutable after training begins. If episodes must be
removed or new demonstrations added, create a new named version so that each
checkpoint can still be linked to the data that produced it.

\begin{checkpointbox}[Interface freeze]
The structure checker passes, every accepted episode has been replayed, camera
frames remain live, actions match operator intent, and the demonstrated pose
range matches the task that will be tested.
\end{checkpointbox}
